\chapter{Background and Related Work}
\label{cha:background}

\section{What Yeast Actually Does}
During fermentation the yeast \emph{Saccharomyces cerevisiae} converts fermentable sugar into
ethanol and carbon dioxide:
\begin{equation}
  \mathrm{C_6H_{12}O_6} \;\longrightarrow\; 2\,\mathrm{C_2H_5OH} + 2\,\mathrm{CO_2}
  \label{eq:fermentation}
\end{equation}
Along the way hundreds of aroma compounds are formed. The most important ones are
\begin{itemize}
  \item \textbf{esters} such as isoamyl acetate (banana) and ethyl hexanoate (apple),
  \item \textbf{higher alcohols} such as isoamyl alcohol (solvent-like),
  \item \textbf{phenols} such as 4-vinylguaiacol (clove, typical for wheat beer).
\end{itemize}

\section{Temperature as a Lever}
Ester formation roughly follows an Arrhenius relation \citep{gerstner2020kinetik}:
\begin{equation}
  k(T) = A \cdot \exp\!\left(-\frac{E_a}{R\,T}\right)
  \label{eq:arrhenius}
\end{equation}
According to Equation~\ref{eq:arrhenius}, a higher temperature should accelerate the formation
considerably -- how strongly, however, is hardly documented for top-fermenting yeasts
\citep{malz2021ester, hefeling2018obergaerig}.

\section{Summary}
Esters dominate the fruity character of top-fermented beers, and their formation depends on
temperature.
